\subsection*{S1.22 Exact reporter-site matching and exploratory natural-variant illustrations}

\paragraph{Question and source population.} We asked whether frozen natural-allele predictions can be interpreted at the same focal variants for which an independent reporter assay measured a regulatory consequence. This analysis revisits public, development-exposed data; it is not a preregistered or blind confirmation. All eleven two-allele reporter loci in the original Figures~4e, 5e and 6c/d of \citet{yang2019asb} were considered before choosing the two main illustrations. The original Supplementary Data~1 workbook \path{41467_2019_9292_MOESM4_ESM.xlsx} supplies mutagenesis primers. The original Source Data workbook \path{41467_2019_9292_MOESM6_ESM.xlsx} supplies the three biological measurements per allele in sheets \path{Figure 4e}, \path{Figure 5e} and \path{Figure 6c}. The latter also contains the SLC39A13 expression measurements associated with the published Figure~6d. Its original label \path{SCL39A13} is explicitly normalised to SLC39A13 in the derived table. No source value or genotype was changed.

\paragraph{Exact identity and allele orientation.} Each pair of mutagenesis primers was matched to the reference and alternate sequences of the saved 101-nt input, checking that its single differing position coincided with focal index~50. Direct and reverse-complement orientations were examined and recorded. Matching uses actual sequences rather than primer names, gene names or pooled read-count coincidences. The initial count-based candidate table is retained to show why those weaker joins were insufficient. For UBE2I only the G mutagenesis primer is supplied; the corresponding C sequence was inferred at its focal G position (primer index~11) from the published C/G alleles, then both were verified against the frozen input and the published hg19 coordinate chr16:1,374,951. That inferred alternative is explicitly flagged. This procedure establishes correspondence to the author's reporter site, not identity of every base in the longer reporter construct.

All eleven sites matched, yielding all 24 available qualified RBP--cell ASB events at those sites and 66 original reporter measurements. Events were retained even when their RBP was not named in the original example, or when the model failed. Twenty-one events are on training chromosomes and three on validation chromosomes; none belongs to the chromosome-8/15 subset. The source protein labels and all three individual conservation permutations are retained. The full cohort of 7,262 ASB events, its fixed 396-event subset and its original statistical endpoints remain unchanged.

\paragraph{Distinct quantities and units.} The source reporter rows give exon or intron inclusion/exclusion in percent, or mCherry RNA expression normalised by eYFP. Their endpoint names, alleles, workbook sheet, row and value column are preserved in \path{ALL66_REPORTER_BIOLOGICAL_MEASUREMENTS.csv}. Each allele has three author-reported biological replicates in HeLa. They are not prediction seeds, patients or new experiments conducted here. Figure~S29 replots every value on its source scale, without converting a splicing percentage into a binding effect. Bars are descriptive arithmetic means with all three measurements shown; no new $p$ value or interval is computed.

BEAPR source preferences and replicate read counts describe allelic binding in K562 or HepG2. Raw read counts are shown for provenance and cannot substitute for BEAPR's bias-adjusted labels or a calibrated binding effect. The model response remains $z(s^+,\Phi)-z(s^-,\Phi)$, in predicted log-total-count units, with the same reference $\Phi$ for both alleles, as in the main-text natural-allele scoring equation. Full-flat and three separate permuted-input responses use fixed Full weights; Sequence is separately trained with unverified matching of training history. Permutation mean is the mean of the three paired responses. Figure~S30 uses one absolute, symmetric asinh scale with linear width 0.02 and retains all 24 events and seven saved arms. Positive/negative symbols decode agreement with the source-preferred allele, not predicted reporter direction. Full and Sequence agree in 12/24 and 11/24 events, respectively; these are descriptive event counts in an exposed, author-selected reporter set, not an inferential performance advantage.

\paragraph{Illustrative cases.} ACVR1B--UPF1 at genomic chr12:51,995,107 A$>$G was selected post hoc because it has one qualified ASB event and an expression reporter at the same site, while Full agrees with the source A preference and Sequence, flat-$\Phi$ and permutation-mean disagree. The responses are $+0.008365$, $-0.002801$, $-0.001236$ and $-0.001097$, respectively. Individual permutations are $-0.004160$, $-0.015574$ and $+0.016443$; permutation~2 also agrees and remains visible. These small differences are not compared numerically with the independent reporter ratios. The source HeLa ratios are A: 1.01, 1.03, 0.96 and G: 1.41, 1.36, 1.46, giving descriptive means 1.00 and 1.41. Directional consistency with UPF1-associated RNA turnover provides a specific follow-up hypothesis, not a new measurement of binding causality.

CERS2 at genomic chr1:150,968,149 T$>$G, RNA A$>$C, is the exact published GWAS-associated rs267738. It was selected for its published allele-specific splicing assay and GEMIN5 knockdown evidence, rather than for Full outperforming Sequence. Both models agree with the GEMIN5 C preference, with C-minus-A responses $+0.024417$ and $+0.073810$. Source exon-inclusion measurements are C: 11.65, 12.89, 13.97 and A: 46.91, 44.21, 46.92 percent. The other qualified same-site RBP, YBX3, remains included: Full predicts the opposite source direction whereas Sequence agrees. The GEMIN5 knockdown result is attributed to the original paper; it is not replotted as newly generated data. GWAS association is distinct from a ClinVar pathogenic classification, and neither implies that this variant is a cancer driver.

\paragraph{Interpretation and display.} Figure~5c shows the post hoc ACVR1B locus, RNA alleles, source binding preference and all seven saved A-minus-G model/input responses on a common linear axis, alongside the distinct HeLa RNA reporter readout. Full/Sequence are separated visually from flat-$\Phi$ and permutation controls. The mean permutation bar and three shape-labelled individual fixed-input responses are all retained; shape IDs 0/1/2 are not biological replicates, refitted seeds or confidence limits. Author reporter points show all three biological measurements per allele and descriptive means. The centre reference PhyloP is printed; both alleles receive the same original 101-nt reference PhyloP vector, and no mutant evolutionary score is inferred. Complete reporter measurements, CERS2 and all associated binding events remain in Figures~S29/S30, with full source sequences/inputs in Data S27. These examples show useful local hypotheses at experimentally examined natural variants. They do not show that CARBON directly predicts exon inclusion, expression, disease risk or RBP-mediated causality. All identities, complete candidate/eligibility census, unsuccessful events, source hashes and selection rules are in Data S27. The ClinVar/TCGA union and its Full 7/13 versus Sequence 10/13 outcome remain unchanged.
